\chapter{Functions, classes, and native data}
\label{ch:interfaces-models-files}

\begin{learningobjectives}
After this chapter, you should be able to:
\begin{itemize}
  \item design a function as a documented public contract;
  \item combine type hints, NumPy-style docstrings, validation, and exceptions;
  \item model a domain object with explicit invariants and state ownership; and
  \item load CSV and JSON without confusing parsing with scientific validation.
\end{itemize}
\end{learningobjectives}

\begin{chapterconnection}
Native values and containers let us represent information. This chapter turns
those representations into interfaces: functions express transformations,
classes protect invariants, exceptions explain rejected inputs, and files carry
representations across process and time boundaries.
\end{chapterconnection}

\section{A function creates a boundary}

A function gives a name to behavior, accepts input through parameters, and
returns a result or raises an exception. The most useful function boundaries
separate one coherent responsibility from its callers.

\begin{definitionbox}{application programming interface}
An application programming interface, or API, is the supported contract through
which one piece of software uses another. A Python function API includes its
import path, signature, parameter meanings, return behavior, exceptions, and
promised side effects. It does not require a network.
\end{definitionbox}

\begin{lstlisting}
def bernoulli_variance(probability: float) -> float:
    """Return the variance of a Bernoulli random variable.

    Parameters
    ----------
    probability : float
        Success probability in the closed interval [0, 1].

    Returns
    -------
    float
        Variance ``p * (1 - p)``.

    Raises
    ------
    TypeError
        If ``probability`` is not numeric.
    ValueError
        If ``probability`` is outside [0, 1].
    """
    if not isinstance(probability, (int, float)):
        raise TypeError("probability must be numeric")
    if not 0.0 <= probability <= 1.0:
        raise ValueError("probability must lie in [0, 1]")
    return probability * (1.0 - probability)
\end{lstlisting}

The annotation helps people, editors, and static checkers. The docstring
explains meaning and failure behavior. Runtime checks enforce selected parts of
the contract. Tests provide independent examples and edge cases. No one layer
replaces the others.

\section{Docstrings explain semantics}

This course uses NumPy-style docstrings for reusable scientific Python. A useful
docstring states the observable behavior, domain meaning, units, shape, ordering,
assumptions, returned representation, and deliberate exceptions. It does not
merely repeat a type annotation.

Consistency matters because readers know where to look, development tools can
render predictable help, and documentation generators can build stable
reference pages. Style is a coordination mechanism.

\section{Exceptions are part of the contract}

Raise an exception when a function cannot honor its promise. Use
\code{TypeError} for an unsupported kind of object and \code{ValueError} for a
value outside the accepted domain. A stable domain exception can help callers
recover without parsing message text.

Library code should generally raise rather than print an error or return an
ambiguous \code{None}. Catch an exception only at a layer that can recover, add
context, translate it for a user, or restore a system invariant.

\begin{warningbox}
An error message is observable output and may enter logs. It should identify the
violated contract without copying API keys, passwords, personal data, or entire
production records.
\end{warningbox}

\section{Parameter kinds communicate calling policy}

Python supports positional-only, positional-or-keyword, variadic positional,
keyword-only, and variadic keyword parameters. The separators \code{/} and
\code{*} make parts of the calling contract explicit.

\begin{lstlisting}
def calibrate(reading, /, *, scale=1.0, offset=0.0):
    return scale * reading + offset
\end{lstlisting}

The name \code{reading} is not promised to callers because the argument is
positional-only. The calibration settings are keyword-only because their names
carry meaning and accidental reversal would remain numerically valid.

\code{*args} and \code{**kwargs} are ordinary tuple and dictionary objects
inside a function. Flexibility is useful for adapters and decorators, but
unrestricted keyword capture can hide spelling errors and weaken tooling.
Lambdas create function objects with expression-only bodies; they are suitable
for short local callables, not for concealing domain logic that deserves a name,
documentation, validation, and tests.

\section{Classes model objects and invariants}

A class combines state with operations that preserve the meaning of that state.
The goal is not to turn every dictionary into a class. A class is valuable when
identity, invariants, behavior, lifecycle, or multiple representations belong
together.

\begin{lstlisting}
from dataclasses import dataclass

@dataclass(frozen=True, slots=True)
class KnowledgeNode:
    identifier: str
    label: str
    domain: str

    def __post_init__(self) -> None:
        if not self.identifier.strip():
            raise ValueError("identifier must be nonblank")
\end{lstlisting}

Immutability can make shared domain records easier to reason about. Mutable
collections can still be owned by a graph or repository whose methods enforce
uniqueness and relationship rules.

\begin{professionalbox}
Represent a scientific object so that invalid states are difficult to create,
but do not invent validation the domain cannot justify. A class can enforce that
a probability lies in $[0,1]$; it cannot establish that the probability was
estimated from an appropriate experiment.
\end{professionalbox}

\section{CSV and JSON are representations}

CSV stores rows and fields in text. It does not carry one universal schema,
type system, missing-value policy, encoding, delimiter, or unit convention.
JSON represents objects, arrays, strings, numbers, Booleans, and null. A parsed
JSON object is not automatically a validated domain object.

A dependable ingestion path separates layers:

\begin{center}
\begin{tikzpicture}[node distance=6mm]
  \node[flowbox] (bytes) {Bytes\\encoding};
  \node[flowbox,right=of bytes] (syntax) {Syntax\\CSV or JSON};
  \node[flowbox,right=of syntax] (schema) {Schema\\fields and types};
  \node[flowbox,right=of schema] (domain) {Domain\\objects and invariants};
  \draw[flowarrow] (bytes) -- (syntax);
  \draw[flowarrow] (syntax) -- (schema);
  \draw[flowarrow] (schema) -- (domain);
\end{tikzpicture}
\end{center}

Use \code{pathlib.Path} rather than string-concatenating separators. Open text
with an explicit encoding. Use \code{newline=""} for Python's CSV reader. For
each row, preserve source location, validate required fields, convert types,
then construct the custom object. Report malformed records with actionable
context without exposing unrelated data.

\begin{checkpointbox}
Design a function that converts one parsed JSON dictionary into a
\code{KnowledgeNode}. Write its signature, NumPy-style contract, two validation
errors, and three tests before implementing its body.
\end{checkpointbox}

\section{Worked example: design the contract before the body}

Suppose we need a function that computes the log-odds

\[
  \operatorname{logit}(p) = \log\!\left(\frac{p}{1-p}\right).
\]

Writing \code{log(p / (1 - p))} is easy. Designing the boundary requires more
thought. The expression is finite only for $0 < p < 1$. The function must decide
whether integers are accepted, whether arrays are in scope, which logarithm is
used, and how boundary values fail.

\begin{workedexample}{turn mathematics into a Python interface}
\textbf{Step 1: state the mathematical domain and codomain.} For this scalar
function, $p\in(0,1)$ and the result is a real-valued natural logarithm.

\textbf{Step 2: state the software representation.} Accept Python \code{int} or
\code{float} values but reject Boolean values, which are integer subclasses in
Python and have no intended role here. Return \code{float}. Arrays will belong
to a separate vectorized interface.

\textbf{Step 3: state failures.} Unsupported runtime types produce
\code{TypeError}; numeric values outside the open interval produce
\code{ValueError}.

\textbf{Step 4: write examples and edge tests.}

\begin{lstlisting}
import math
import pytest

assert math.isclose(logit(0.5), 0.0)
assert logit(0.8) > 0.0
with pytest.raises(ValueError):
    logit(1.0)
with pytest.raises(TypeError):
    logit(True)
\end{lstlisting}

\textbf{Step 5: implement the smallest body that honors the contract.}

\begin{lstlisting}
from math import log

def logit(probability: float) -> float:
    """Return the natural log-odds of a scalar probability."""
    if isinstance(probability, bool) or not isinstance(
        probability, (int, float)
    ):
        raise TypeError("probability must be a real scalar")
    if not 0.0 < probability < 1.0:
        raise ValueError("probability must lie strictly between 0 and 1")
    return log(probability / (1.0 - probability))
\end{lstlisting}
\end{workedexample}

The tests did not emerge after the implementation. They helped define what the
function means. A later NumPy implementation may accept arrays and infinities at
the boundaries; that is a different public contract, not a faster body for the
same one.

\section{From records to objects with behavior}

A dictionary can store a graph node, but any caller can omit a field, replace an
identifier with whitespace, or mutate a label without going through validation.
A class centralizes construction and behavior.

Consider a \code{KnowledgeGraph} that owns nodes and adjacency relationships.
Its public methods might be \code{add\_node}, \code{add\_edge},
\code{neighbors}, and \code{shortest\_path}. Internal dictionaries are private
because callers who mutate them directly can violate uniqueness or create an
edge to a nonexistent node.

This is \term{encapsulation}: state is accessed through operations that preserve
the object's invariants. It is not secrecy. Python callers may still inspect
attributes; the API communicates which path is supported.

\begin{misconceptionbox}
\textbf{Misconception:} object-oriented design means every noun needs a class.

\textbf{Correction:} functions and native containers are often clearer. Use a
class when state, behavior, invariants, identity, or lifecycle form one coherent
abstraction. Avoid classes that only rename a dictionary without improving the
contract.
\end{misconceptionbox}

\section{A complete ingestion boundary}

Loading a knowledge graph from files involves four different failure layers:

\begin{enumerate}
  \item \textbf{I/O failure:} the path does not exist or cannot be read;
  \item \textbf{syntax failure:} bytes are not valid text, CSV, or JSON;
  \item \textbf{schema failure:} a required field is missing or has the wrong
  representation; and
  \item \textbf{domain failure:} an edge references an unknown node or violates
  a graph invariant.
\end{enumerate}

Keeping these layers distinct makes errors actionable. ``Could not load data''
hides whether a student mistyped a path or discovered a corrupted relationship.
A high-quality loader adds source path and record number while preserving the
original exception through chaining.

\begin{lstlisting}
try:
    node = node_from_mapping(row)
except (KeyError, TypeError, ValueError) as error:
    raise ValueError(
        f"invalid node record at {path.name}:{line_number}"
    ) from error
\end{lstlisting}

The public message supplies location without dumping the full row. The chained
exception preserves technical cause for debugging.

\conceptcheck{Why should a loader validate that both endpoints exist before
adding an edge, even if a later graph algorithm would eventually fail? Compare
the quality of the two failure locations.}

\section{Exercises}

\subsection*{Foundational}
\begin{enumerate}
  \item For a function you used recently, list its import path, parameter
  meanings, return behavior, exceptions, and side effects. Which are documented?
  \item Explain why a type hint does not perform runtime validation by itself.
  \item Contrast \code{*args} at a function definition with \code{*values} at a
  function call.
\end{enumerate}

\subsection*{Applied}
\begin{enumerate}
  \item Complete the logit docstring using NumPy style, including domain,
  return meaning, and exceptions. Add tests at values close to 0 and 1.
  \item Design an immutable \code{ChemicalSpecies} dataclass with formula,
  molar mass, and phase. State what can be validated syntactically and what
  requires an external scientific source.
  \item Create a three-row CSV fixture containing one schema error and one
  domain error. Specify the exact message and chained exception expected for
  each.
\end{enumerate}

\subsection*{Design}
Design the public API of a small graph, trajectory, molecule, queue, or
experiment object. Separate public operations from private representation and
state at least three invariants each operation must preserve.

\begin{chapterreview}
\textbf{Key ideas.} A function is a boundary, not merely reusable syntax.
Docstrings explain semantics, types communicate representation, validation
enforces selected rules, exceptions communicate failure, and tests supply
independent examples. Classes are justified by coherent state and invariants.
File parsing precedes schema and domain validation.

\textbf{Vocabulary.} function; signature; parameter; return value; API;
docstring; type hint; exception; keyword-only; lambda; class; dataclass;
invariant; encapsulation; CSV; JSON; exception chaining.

\textbf{Explain aloud.} Why can two functions compute the same formula but have
different APIs?
\end{chapterreview}

\begin{notebooklink}
Lecture 2 develops functions, classes, lambdas and variadic calls, then reads
real CSV and JSON relationships into a custom knowledge graph using only the
standard library.
\end{notebooklink}

\section{Takeaway}

Functions and classes turn assumptions into named interfaces. Files supply
representations, not meaning. Professional code makes the transition from one
to the other explicit, documented, testable, and diagnosable.
